\documentclass[11pt]{article}

\usepackage{acl}

\usepackage{times}
\usepackage{latexsym}
\usepackage[T1]{fontenc}
\usepackage[utf8]{inputenc}
\usepackage{microtype}
\usepackage{inconsolata}
\usepackage{graphicx}
\usepackage{amsmath}
\usepackage{booktabs}
\usepackage{hyperref}
\usepackage{enumitem}

\title{Stability by Design: Optimizing Prompt Properties in Small Language Models via Sensitivity Metrics}

\author{
<Full Name 1> (ID: 209127612) \\
\texttt{<email1>}
\And
Sharon Levy (ID: 315140152) \\
\texttt{sharon14@mail.tau.ac.il}
\AND
<Full Name 3> (ID: 315285759) \\
\texttt{<email3>}
\And
<Full Name 4> (ID: 314948811) \\
\texttt{<email4>}
\\[0.5em]
\multicolumn{2}{c}{\textbf{Moodle Group Number: <GROUP NUMBER>}}
}

\begin{document}
\maketitle

\begin{abstract}
UPDATE LATER
\end{abstract}

\section{Introduction}

Small language models (SLMs) are becoming an attractive option for practical NLP use because they are cheaper to run, easier to host locally, and more suitable for privacy-sensitive or resource-constrained settings than frontier-scale models. However, these advantages come with an important limitation. Model behavior can depend strongly on the way a task is phrased. In practice, small changes to a prompt, such as synonym substitutions, paraphrases, added conversational phrasing, or different output formats, may change the model's prediction and, as a result, its measured task performance.

This issue matters both scientifically and practically. From a scientific point of view, prompt instability makes evaluation harder to interpret. If two prompts express the same task intent but lead to different outputs, then part of the observed performance may reflect wording choices rather than the model's actual ability. From a practical point of view, instability reduces reliability. In real applications, prompts are often reformulated by users, developers, or other components in a pipeline. Ideally, these reformulations should not cause large and unpredictable changes in model behavior. This concern is especially relevant for SLMs, where limited capacity and weaker instruction-following may increase the effect of superficial prompt variation.

At the same time, sensitivity is not meaningful on its own. A model that gives the same incorrect answer across all prompt variants may appear stable, even though it is not robust in any useful sense. For this reason, prompt robustness should be studied together with baseline task competence. In this project, we study prompt robustness in SLMs through the lens of \textbf{input sensitivity}, while also checking whether a model can perform the underlying task well enough for sensitivity scores to be meaningful.

Following \citet{lu-etal-2024-prompts}, we define prompt sensitivity as the extent to which a model's prediction changes when the original prompt is replaced with lightly perturbed variants that preserve its meaning. We adapt this framework to small, locally runnable models and use a controlled perturbation pipeline that generates semantic-preserving variants through synonym replacement and paraphrasing. Our goal is not only to measure instability, but also to ask whether certain prompt properties systematically affect it.

We focus on three prompt properties inspired by prior work on prompt design \citep{long-etal-2025-makes}. The first is \textbf{metacognition}, implemented through self-check or self-verification cues. The second is \textbf{structure}, implemented through explicit output constraints, in our case a fixed JSON response format. The third is \textbf{politeness}, implemented through courteous conversational phrasing. These conditions are compared against a standard zero-shot control prompt.

\textbf{Research Question:} Do certain prompt properties (metacognition, structure, politeness) make language models more stable when inputs undergo semantic-preserving perturbations?

To answer this question, we use an inference-only evaluation framework. For each prompt style, we run the same model on the original prompt and on several semantic-preserving perturbations, and we measure sensitivity using the variation ratio,
\[
s = 1 - \frac{f_m}{N+1},
\]
where $f_m$ is the frequency of the modal prediction across the original prompt and its $N$ perturbed variants. Lower values indicate greater consistency across perturbations. We evaluate this framework on two active benchmarks with different task structures: QASC, an 8-way multiple-choice science question answering task, and CoLA, a binary grammatical acceptability task. We test a suite of locally runnable models that includes both encoder-decoder and decoder-only architectures, as well as base and instruction-tuned variants.

Our contributions are as follows:
\begin{itemize}
    \item We formulate prompt robustness in SLMs as a measurable sensitivity problem, while emphasizing that sensitivity is only interpretable when the model shows non-trivial baseline task competence.
    \item We adapt a sensitivity-based evaluation framework to small, locally runnable models using a controlled perturbation pipeline based on semantic-preserving synonym replacement and paraphrasing.
    \item We analyze whether three prompt properties, metacognition, structure, and politeness, affect both sensitivity and task performance across two different benchmark settings.
\end{itemize}

\section{Related Work}

Our project connects two lines of research: work on \textbf{prompt design and prompt attributes}, and work on \textbf{prompt sensitivity and input stability}.

The first line of work asks what makes a prompt effective. Recent research has moved beyond treating prompts as isolated heuristics and instead describes them in terms of reusable and interpretable attributes. In particular, \citet{long-etal-2025-makes} propose a taxonomy of prompt attributes that organizes prompt design choices into dimensions such as metacognitive guidance, structural constraints, and stylistic framing. This perspective is central to our project because it makes prompt design easier to study systematically rather than only through trial and error. We build directly on this view by selecting three attributes, metacognition, structure, and politeness, and operationalizing them in a controlled experimental setup.

The second line of work studies prompt quality through robustness to input variation. \citet{lu-etal-2024-prompts} analyze prompts in terms of \textbf{input stability}, showing that prompts differ substantially in their sensitivity to minor perturbations and that higher sensitivity often accompanies worse downstream performance. This work motivates our project in two ways. First, it suggests that prompt evaluation should consider not only accuracy, but also whether predictions remain stable under semantically equivalent reformulations. Second, it provides a concrete sensitivity metric, variation ratio, that can be used as a diagnostic signal during prompt evaluation.

Our work builds on both lines of research, but it also differs from them in several important ways. Compared with prior work on prompt attributes, we are not only interested in whether a prompt improves average performance, but also in whether it changes robustness under controlled perturbations. Compared with prior work on sensitivity, we focus specifically on \textbf{small, locally runnable models}, including both base and instruction-tuned variants, where prompt sensitivity may interact strongly with limited instruction-following ability and with the model's underlying task competence. In this setting, low sensitivity is not always desirable if it simply reflects consistently wrong behavior. For this reason, our framework separates baseline task competence from sensitivity analysis by checking out-of-the-box (OOTB) accuracy before interpreting robustness scores.

The gap we address lies in the connection between \textbf{how prompts differ in form} and \textbf{how prompts differ in robustness}. Work on prompt attributes explains how prompts can vary in structure, intention, and style. Work on sensitivity shows that prompts can vary in stability under perturbation. What remains less explored is the relation between these two questions in small-model settings: which prompt properties make prompts more or less stable, when greater stability aligns with better task performance, and when stability can become a misleading signal. Our project addresses this gap by combining an attribute-based view of prompting with a sensitivity-based evaluation framework.

\section{Method}

\subsection{Task Definition}

We study prompt robustness as a conditional prediction problem under semantic-preserving reformulation. Let $x$ denote an input instance, let $y$ denote its gold label, and let $p$ denote a prompt style. A model receives a prompted version of the input, produces a prediction $\hat{y}$, and is then evaluated both on the original input and on lightly perturbed variants of that same input.

Our setting includes two task types. In QASC, the input is a multiple-choice science question with eight answer options, and the target output is a single answer letter from A to H. In CoLA, the input is a sentence and the target output is a binary grammaticality judgment, represented as Yes or No. In both cases, the task content remains fixed across prompt styles; only the framing of the instruction changes.

For each instance, we compare four prompt styles. The \textbf{control} condition is a standard zero-shot instruction. The \textbf{metacognition} condition adds self-check and verification cues. The \textbf{structure} condition constrains the model to respond in a fixed JSON format that separates a brief explanation from a final answer field. The \textbf{politeness} condition adds courteous conversational phrasing. The goal is not to optimize prompts for a single benchmark, but to test whether these prompt properties change the stability of model predictions under controlled reformulations of the same underlying input.

\subsection{Approach}

Our framework is inference-only. We do not fine-tune any model parameters. Instead, we evaluate how the same pretrained model behaves when the same task is presented through different prompt styles and through semantic-preserving input perturbations.

For each model-dataset pair, we first run an OOTB accuracy check using the control prompt. This step is meant to verify that the model can perform the task at a level that makes sensitivity analysis meaningful. Only after this check do we run the full sensitivity experiment.

In the main experiment, we sample a set of input instances and evaluate each prompt style separately. For every sampled instance, we construct one original prompt and $N$ perturbed variants. The model is run on all $N+1$ versions, and the resulting predictions are parsed into task-level labels. For QASC, this means extracting one answer letter. For CoLA, this means extracting a Yes or No judgment. In the structure condition, the parser reads the final answer field from the JSON output; in the other conditions, it extracts the relevant label from plain text output.

We use two perturbation strategies. The default strategy is synonym-based perturbation. The input text is tokenized and part-of-speech tagged, and candidate replacement words are restricted to content words such as nouns, verbs, adjectives, and adverbs. Very short words, stopwords, question words, and answer labels are excluded. Candidate synonyms are taken from WordNet, with preference for the most common senses in order to reduce semantic drift. When synonym replacement cannot produce enough distinct variants, the pipeline uses lightweight fallback transformations such as minor article changes, filler prefixes, or punctuation variation.

The second strategy is paraphrasing. Here, a separate small model is used to rewrite the input while preserving its meaning. The paraphrase generator cycles through multiple rewrite instructions and uses stochastic decoding to produce diverse variants. Duplicate paraphrases are removed, and any shortfall is again completed by the synonym-based fallback. In both perturbation modes, the final goal is the same: produce multiple variants that change surface form while preserving task intent as much as possible.

For each instance and prompt style, we compute sensitivity using the variation ratio over the model's parsed predictions. We then average this quantity across all sampled instances. In parallel, we also track task accuracy so that low sensitivity can be interpreted together with baseline competence rather than in isolation.

\section{Experimental Setup}

\subsection{Data}

We evaluate the framework on two benchmarks with different task structures.

The first is QASC, an 8-way multiple-choice science question answering benchmark. Each example contains a question and eight candidate answers. In early pilot runs, including the dataset's supporting facts made the task too easy for stronger instruction-tuned models, producing near-ceiling behavior and leaving little room for meaningful sensitivity comparisons. For this reason, the main setup uses QASC without fact injection. The model therefore sees the question and answer choices, but not the additional supporting facts. This makes the benchmark more suitable for studying whether prompt style changes both robustness and performance under a non-trivial level of difficulty.

The second benchmark is CoLA, a binary grammatical acceptability task. Each example consists of a single sentence, and the model must decide whether it is grammatically acceptable. CoLA complements QASC in two useful ways. First, it has a much simpler output space. Second, it evaluates a different kind of linguistic competence, namely grammatical judgment rather than multiple-choice reasoning.

All experiments use the validation split of each benchmark. For the main sensitivity experiment, the framework samples up to 500 examples per model-dataset pair using a fixed seed. For each sampled example, it generates 10 perturbations in addition to the original input. The OOTB accuracy check is run on a separate sample of 100 instances before the sensitivity experiment begins.

\subsection{Models}

We evaluate a set of locally runnable models chosen to vary along two dimensions: architecture and instruction tuning. The model suite includes both encoder-decoder and decoder-only systems, as well as both base and instruction-tuned variants.

Our encoder-decoder models are Flan-T5-Base and Flan-T5-Large. These models are useful reference points because they are explicitly instruction-tuned and relatively strong on zero-shot text-to-text tasks. Our decoder-only base models are Pythia-410M and Llama-3.2-1B. These models let us examine whether prompt sensitivity behaves differently when instruction-following ability is limited. Our instruction-tuned decoder-only models are Llama-3.2-1B-Instruct and Phi-3-Mini-4K-Instruct. These provide a comparison point within the decoder-only family between base and instruction-tuned behavior.

All models are loaded directly from HuggingFace and evaluated without any parameter updates. For causal models, padding is set on the left and the end-of-sequence token is used as a padding token when necessary. For instruction-tuned chat models, the framework automatically applies the tokenizer's chat template before generation. This keeps model-specific formatting separate from the core evaluation loop and ensures that each model is used in a way that is consistent with its intended interface.

\subsection{Training}

Because our study focuses on prompt sensitivity rather than adaptation or fine-tuning, there is no training stage in the usual sense. All experiments are performed with frozen pretrained models and deterministic inference.

Generation is run with greedy decoding so that differences in output are driven by prompt and input changes rather than by sampling noise in the evaluated model itself. The maximum number of generated tokens is task- and style-dependent. Plain answer formats use shorter generation limits, while the structure condition allows longer outputs because the model must produce a valid JSON object containing both a brief explanation and a final answer field.

The implementation is designed for reproducibility. Random seeds are fixed across Python, NumPy, PyTorch, and the Transformers library. The framework defines three experiment seeds, 105, 2266, and 86379, and each run is fully determined by its seed and configuration. The perturbation pipeline also uses seed control at the level of individual examples so that the same configuration reproduces the same perturbation set across runs. Models are executed on the best available device, including CUDA, Apple MPS, or CPU, and half precision is used for selected decoder-only models when supported.

\section{Evaluation}

We evaluate each model-prompt combination along two axes: task performance and sensitivity.

The first metric is \textbf{OOTB accuracy}. Before running sensitivity experiments, we measure the model's accuracy on unperturbed inputs using the control prompt. This serves as a sanity check. If a model performs at or near random guessing, then low sensitivity is not informative, because the model may simply be repeating the same wrong answer. In our implementation, this check is dataset-specific. For QASC, the random baseline is 12.5\% and the valid-signal threshold is set to 40\%. For CoLA, the random baseline is 50\% and the valid-signal threshold is set to 60\%.

The second metric is \textbf{variation ratio}, which measures the dispersion of predictions across the original input and its perturbations:
\[
s = 1 - \frac{f_m}{N+1},
\]
where $f_m$ is the frequency of the modal prediction across the original prompt and its $N$ perturbed variants. Lower values indicate greater consistency, while higher values indicate greater sensitivity to reformulation. Importantly, we compute this score over parsed task labels rather than raw generated strings. This is especially important for the structure condition, where the explanatory text may vary even when the final answer remains unchanged.

In addition to average variation ratio, we also report \textbf{style-specific accuracy} on the original, unperturbed inputs used in the main experiment. This lets us compare prompt styles not only in terms of stability, but also in terms of whether increased consistency coincides with better or worse task performance. Together, these metrics allow us to distinguish several cases: prompts that are both accurate and stable, prompts that are accurate but unstable, and prompts that appear stable only because the model gives the same incorrect answer repeatedly.


\section{Results}

\begin{table}[h]
\centering
\begin{tabular}{l c}
\toprule
Model & Score \\
\midrule
Baseline & XX \\
Ours & XX \\
\bottomrule
\end{tabular}
\caption{Results}
\end{table}

\section{Discussion}

<Interpret results>

\section{Limitations}

<Constraints, weaknesses>

\section{Conclusion}

<Summary>

\section*{AI Usage Disclosure}

We used AI tools for <X>. We did not use AI for generating results or data.
All decisions and evaluations were performed by us.

\bibliography{custom}

\end{document}